\documentclass{article}

% Recommended, but optional, packages for figures and better typesetting:
\usepackage{microtype}
\usepackage{graphicx}
\usepackage{subfigure}
\usepackage{booktabs} % for professional tables

% hyperref makes hyperlinks in the resulting PDF.
% If your build breaks (sometimes happens with hyperref and
% latex), clean out all the auxiliary files from your PDF
% directory and try again.
\usepackage{hyperref}


% Attempt to make hyperref and algorithmic work together better:
\newcommand{\theHalgorithm}{\arabic{algorithm}}

% Use the following line for the initial blind version submitted for review:
\usepackage[accepted]{icml2026}

% If accepted, change the following line for the camera-ready version:
% \usepackage[accepted]{icml2026}

% For theorems and such
\usepackage{amsmath}
\usepackage{amssymb}
\usepackage{mathtools}
\usepackage{amsthm}

% if you use cleveref..
\usepackage[capitalize,noabbrev]{cleveref}

%%%%%%%%%%%%%%%%%%%%%%%%%%%%%%%%
% THEOREMS
%%%%%%%%%%%%%%%%%%%%%%%%%%%%%%%%
\theoremstyle{plain}
\newtheorem{theorem}{Theorem}[section]
\newtheorem{proposition}[theorem]{Proposition}
\newtheorem{lemma}[theorem]{Lemma}
\newtheorem{corollary}[theorem]{Corollary}
\theoremstyle{definition}
\newtheorem{definition}[theorem]{Definition}
\newtheorem{assumption}[theorem]{Assumption}
\theoremstyle{remark}
\newtheorem{remark}[theorem]{Remark}

% Todonotes is useful during development; simply uncomment the next line
%    and comment out the line below the next line to turn off comments
%\usepackage[disable,textsize=tiny]{todonotes}
\usepackage[textsize=tiny]{todonotes}


% The \icmltitle you define below is probably too long as a header.
% Therefore, a short form for the running title is supplied here:
\icmltitlerunning{Activation Overlap Search}

\begin{document}

\twocolumn[
\icmltitle{Activation Overlap Search: Data-Driven Estimation of Model Merging Coefficients}

% It is OKAY to include author information, even for blind
% submissions: the style file will automatically remove it for you
% unless you've provided the [accepted] option to the icml2026
% package.

% List of affiliations: The first argument should be a (short)
% identifier you will use later to specify author affiliations
% Academic affiliations should list Department, University, City, Region, Country
% Industry affiliations should list Company, City, Region, Country

% You can specify symbols, otherwise they are numbered in order.
% Ideally, you should not use this facility. Affiliations will be numbered
% automatically, and symbols will be ignored.

\icmlsetsymbol{equal}{*}

\begin{icmlauthorlist}
\icmlauthor{Anonymous Authors}{yyy}
\end{icmlauthorlist}

\icmlaffiliation{yyy}{Department of Empirical Machine Learning, Institute for Large-Scale Experimentation, Computron, Earth}

\icmlcorrespondingauthor{}{}

% You may provide any keywords that you
% find helpful for describing your paper; these are used to populate
% the "keywords" metadata in the PDF file.
\icmlkeywords{Model Merging, Deep Learning, Representation Similarity, Empirical Methods}

\vskip 0.3in
]

% this must go after the closing bracket ] following \twocolumn[ ...

% This command actually creates the footnote in the first column
% listing the affiliations and the copyright notice.
% The command takes one argument, which is text to display at the start of the footnote.
% The \icmlEqualContribution command is standard text for equal contribution.
% Remove it if you have shared first authorship.
\printAffiliationsAndNotice{}  % leave blank if no need to mention equal contribution
% \printAffiliationsAndNotice{\icmlEqualContribution} % otherwise use the standard text.

\begin{abstract}
Model merging aims to combine multiple fine-tuned models into a single multitask model, but its success often hinges on finding the right coefficients to balance the models' parameters. The standard approach---a costly grid search over coefficients evaluated on a validation set---undermines the computational promise of merging. We introduce Activation Overlap Search (AOS), a data-driven, training-free method to estimate these coefficients. AOS operates on the principle that the optimal merge coefficient is the one that maximally preserves the functional behavior of a base model, which we measure by comparing internal activation manifolds. Through a massive empirical study spanning over 2,160 unique configurations across three datasets, four similarity metrics, and numerous hyperparameters, we demonstrate that AOS with Centered Kernel Alignment (CKA) is a near-perfect proxy for validation accuracy. Our method exactly identifies the oracle peak accuracy of 58.43\%, a feat unachievable by naive Weight Averaging, which catastrophically fails at 34.58\%. Furthermore, AOS is remarkably data-efficient, nearing oracle performance with as few as 16 data samples. Our work provides overwhelming empirical evidence that functional similarity in activation space is a powerful and practical tool for navigating the model merging landscape.
\end{abstract}

\section{Introduction}

The paradigm of fine-tuning large pre-trained models on specific downstream tasks has become a cornerstone of modern machine learning. This often results in a proliferation of specialist models, each adept at a single task but collectively consuming significant storage and computational resources. Model merging offers an elegant solution: combining the capabilities of these specialist models into a single, powerful generalist, ideally without the exorbitant cost of retraining from scratch.

A common and effective merging strategy is a simple linear interpolation in weight space, such as the task arithmetic framework \citep{ilharco2022editing}. Here, the parameters of a merged model $\theta_{merged}$ are formed by combining a base model $\theta_{base}$ and a task-specific model $\theta_{task}$ via a scaling coefficient $\lambda$: $\theta_{merged} = (1-\lambda)\theta_{base} + \lambda\theta_{task}$. While simple, this approach introduces a critical hyperparameter, $\lambda$. The performance of the merged model is extremely sensitive to this coefficient. A poorly chosen $\lambda$ can lead to a model that performs worse than either of its parents, a phenomenon we demonstrate empirically.

The standard method for finding an optimal $\lambda$ involves a brute-force grid search: for a range of $\lambda$ values, one creates the merged model and evaluates its performance on a held-out validation set. This "line search" is computationally expensive, requiring numerous forward passes over a potentially large dataset for each pair of models to be merged. For large language models or when merging many tasks, this cost becomes prohibitive, defeating the purpose of a "training-free" merging procedure.

We argue that true progress in machine learning comes from exhaustive empirical validation. To that end, we questioned whether this expensive, evaluation-driven search for $\lambda$ could be replaced with a cheaper, data-driven proxy. Our core hypothesis is that the optimal merge coefficient $\lambda^*$ should produce a merged model whose internal function is as close as possible to one of the parent models. Instead of measuring functional similarity via expensive downstream task evaluation, we propose to measure it directly by comparing the models' internal activation manifolds on a small set of probe data.

In this paper, we introduce **Activation Overlap Search (AOS)**, a simple, general, and computationally cheap method for estimating model merging coefficients. AOS iterates through candidate $\lambda$ values and, for each, computes the similarity between the activations of the resulting merged model and a reference model. The $\lambda$ that maximizes this activation similarity is chosen as the optimal coefficient. This process requires no labels, no backpropagation, and as our results show, remarkably few data samples.

To validate AOS, we conducted a massive-scale empirical investigation. We systematically explored the efficacy of AOS across three vision datasets (CIFAR-10, SVHN, FashionMNIST), four distinct activation similarity metrics (CKA, MSE, Cosine, MMD), six different sample sizes for activation computation (from N=8 to N=256), three different network layers, and five independent random seeds for every configuration. This amounts to over 2,160 unique experimental runs, providing a robust and comprehensive picture of our method's performance.

Our findings are decisive. We show that AOS using Centered Kernel Alignment (CKA) as the similarity metric is a nearly perfect proxy for the oracle grid search, consistently identifying a $\lambda$ that yields peak performance. Furthermore, we identify and solve a critical but often overlooked issue with merging models containing Batch Normalization layers, proposing a simple fix that stabilizes the merging process and prevents catastrophic performance collapse.

\section{Related Work}

Model merging is a rapidly developing field aimed at combining the knowledge of multiple neural networks. Early approaches focused on simple averaging of model weights, often referred to as Weight Averaging (WA), which works surprisingly well for models trained from the same initialization but often fails for models in different basins of the loss landscape.

**Task Arithmetic.** \citet{ilharco2022editing} introduced the concept of task vectors, representing the difference between fine-tuned and pre-trained model weights ($\tau = \theta_{task} - \theta_{base}$). Merging is then performed by adding a scaled task vector to the base model: $\theta_{merged} = \theta_{base} + \lambda\tau$. This formulation makes the role of the scaling coefficient $\lambda$ explicit and is the primary setting we address.

**Permutation-based Alignment.** A significant challenge in merging is that functionally similar neurons can occupy different indices in the weight tensors of independently trained models. Several methods seek to find permutations to align these neurons before averaging. Git Re-Basin \citep{ainsworth2022git} finds such permutations by analyzing weight correlations, while SyMerge \citep{yadav2023symerge} uses a small amount of data to align models. OrthoMerge \citep{pena2024orthomerge} generalizes this by finding optimal orthogonal transformations to align weight spaces. While powerful, these methods add significant computational complexity compared to simple linear interpolation. AOS, in contrast, focuses on finding the best interpolation coefficient for a given alignment (or lack thereof).

**Sparsity and Sign Consensus.** Other lines of work focus on the weight values themselves. TIES-Merging \citep{yadav2023ties} addresses the problem of sign disagreement in task vectors by resetting conflicting weights and merging only the most significant parameter changes. DARE \citep{yu2023dare} takes this further by dropping a significant fraction of task vector weights and rescaling the remainder, achieving strong results with high sparsity. These methods are complementary to ours; they modify the task vector $\tau$, while AOS finds the optimal scaling $\lambda$ for it.

**Coefficient Learning.** Few works have focused directly on automating the discovery of merging coefficients. TCAC \citep{pan2024tcac} proposes learning task-complementary attention condensers to derive merging weights, but this involves a training phase. AOS is explicitly training-free. Our work is the first, to our knowledge, to conduct a large-scale empirical study of using activation similarity as a direct, training-free proxy for finding scalar merging coefficients.

\section{Methodology}

The core objective of Activation Overlap Search (AOS) is to efficiently estimate the optimal scalar coefficient $\lambda^*$ for merging two models, $\theta_A$ and $\theta_B$, without performing a full evaluation on a validation set. We focus on the linear interpolation or task arithmetic setting, where the merged model's parameters $\theta_\lambda$ are defined as:
\begin{equation}
\theta_\lambda = (1-\lambda)\theta_A + \lambda\theta_B
\end{equation}
The central hypothesis of AOS is that the optimal $\lambda^*$ that maximizes downstream performance is approximately equal to the $\hat{\lambda}$ that maximizes the functional similarity between the merged model $f_\lambda$ and a reference model (e.g., $f_A$). We propose to measure this functional similarity by comparing their internal representations.

\subsection{The AOS Algorithm}

The AOS procedure is as follows:
\begin{enumerate}
    \item \textbf{Select Data:} Choose a small, unlabeled set of $N$ data points, $X = \{x_1, \dots, x_N\}$.
    \item \textbf{Define Search Space:} Define a set of candidate coefficients, e.g., $\Lambda = \{0.0, 0.1, \dots, 1.0\}$.
    \item \textbf{Extract Reference Activations:} Perform a forward pass with the reference model (e.g., model A) on the data $X$ and extract the activations from a chosen layer $l$. This results in an activation matrix $M_A^{(l)} \in \mathbb{R}^{N \times D_l}$, where $D_l$ is the dimension of the flattened activation at layer $l$.
    \item \textbf{Search for Optimal Overlap:} For each candidate $\lambda \in \Lambda$:
    \begin{enumerate}
        \item Construct the merged model $\theta_\lambda$.
        \item Extract its activations at layer $l$ on data $X$ to get $M_\lambda^{(l)}$.
        \item Compute a similarity score $S(\lambda) = \text{Similarity}(M_A^{(l)}, M_\lambda^{(l)})$.
    \end{enumerate}
    \item \textbf{Estimate Coefficient:} The estimated optimal coefficient $\hat{\lambda}$ is the one that maximizes the similarity score:
    \begin{equation}
    \hat{\lambda} = \arg\max_{\lambda \in \Lambda} S(\lambda)
    \end{equation}
\end{enumerate}
The resulting $\hat{\lambda}$ is then used to construct the final merged model. This entire process is training-free and typically takes only a few seconds to minutes, depending on model size and the number of data points $N$.

\subsection{Similarity Metrics}
The choice of similarity metric is crucial. We investigated four widely-used metrics for comparing representations. For all metrics, we define $S(\lambda)$ such that higher values indicate greater similarity.

\textbf{Centered Kernel Alignment (CKA).} CKA \citep{kornblith2019similarity} is a robust metric of representation similarity, invariant to orthogonal transformations and isotropic scaling. Given two activation matrices $M_1, M_2 \in \mathbb{R}^{N \times D}$, we first form the Gram matrices $K_1 = M_1 M_1^T$ and $K_2 = M_2 M_2^T$. CKA is defined as the normalized Hilbert-Schmidt Independence Criterion (HSIC) of these kernels:
\begin{equation}
\text{CKA}(M_1, M_2) = \frac{\text{HSIC}(K_1, K_2)}{\sqrt{\text{HSIC}(K_1, K_1) \text{HSIC}(K_2, K_2)}}
\end{equation}
where $\text{HSIC}(K_1, K_2) = \frac{1}{(N-1)^2}\text{Tr}(HK_1HK_2)$ and $H = I - \frac{1}{N}\mathbf{1}\mathbf{1}^T$ is the centering matrix.

\textbf{Mean Squared Error (MSE).} A direct way to measure similarity is the distance between activation vectors. We use the negative MSE to frame it as a maximization problem:
\begin{equation}
S(\lambda) = - \frac{1}{N D_l} \sum_{i=1}^N \| f_A^{(l)}(x_i) - f_\lambda^{(l)}(x_i) \|_2^2
\end{equation}

\textbf{Cosine Similarity.} To abstract away from the magnitude of activations, we can use the average cosine similarity between activation vectors:
\begin{equation}
S(\lambda) = \frac{1}{N} \sum_{i=1}^N \frac{f_A^{(l)}(x_i) \cdot f_\lambda^{(l)}(x_i)}{\| f_A^{(l)}(x_i) \| \| f_\lambda^{(l)}(x_i) \|}
\end{equation}

\textbf{Maximum Mean Discrepancy (MMD).} MMD measures the distance between the distributions of activations. Using a radial basis function (RBF) kernel $k(u,v) = \exp(-\|u-v\|^2 / 2\sigma^2)$, the squared MMD is computed. We again use the negative to create a similarity score:
\begin{align}
\text{MMD}^2 = &\frac{1}{N^2} \sum_{i,j} k(f_A^{(l)}(x_i), f_A^{(l)}(x_j)) \nonumber \\
                - &\frac{2}{N^2} \sum_{i,j} k(f_A^{(l)}(x_i), f_\lambda^{(l)}(x_j)) \nonumber \\
                + &\frac{1}{N^2} \sum_{i,j} k(f_\lambda^{(l)}(x_i), f_\lambda^{(l)}(x_j))
\end{align}

\subsection{Handling Batch Normalization}
A crucial and overlooked aspect of model merging is the handling of Batch Normalization (BN) layers \citep{ioffe2015batch}. BN layers contain not only learnable affine parameters ($\gamma$, $\beta$), which are part of $\theta$ and are interpolated, but also non-trainable running statistics (\texttt{running\_mean}, \texttt{running\_var}) which are essential for inference. Naively interpolating only the weights while leaving the running statistics untouched leads to a severe mismatch between activations and normalization statistics.
Our empirical investigation revealed that this mismatch causes activation variances to collapse, leading to numerical instability (NaNs) and catastrophic performance drops, particularly for $\lambda > 0.3$. To remedy this, we propose to interpolate the BN running statistics alongside the weights:
\begin{align}
\text{mean}_\lambda &= (1-\lambda)\text{mean}_A + \lambda\text{mean}_B \\
\text{var}_\lambda &= (1-\lambda)\text{var}_A + \lambda\text{var}_B
\end{align}
This simple fix proved to be absolutely critical for stabilizing the merge process and enabling a fair and complete search over the $\lambda \in [0, 1]$ space.

\subsection{Multidimensional AOS (M-AOS)}
While single-coefficient searching is effective, complex multi-task environments often benefit from task-specific scaling factors. We introduce Multidimensional Activation Overlap Search (M-AOS) to scale our framework to multi-task regimes where each expert model $t$ has its own task coefficient $\lambda_t$:
\begin{equation}
\theta_{\vec{\lambda}} = \theta_{\text{base}} + \sum_{t=1}^T \lambda_t (\theta_t - \theta_{\text{base}})
\end{equation}
M-AOS finds the optimal task-specific parameter vector $\vec{\lambda}^* = (\lambda_1, \dots, \lambda_T)$ by maximizing the average activation similarity across all tasks simultaneously:
\begin{equation}
\vec{\lambda}^* = \arg\max_{\vec{\lambda}} \frac{1}{T} \sum_{t=1}^T \text{Similarity}(M_{\vec{\lambda}, t}^{(l)}, M_{t}^{(l)})
\end{equation}
where $M_{\vec{\lambda}, t}^{(l)}$ are the activations of the merged model $\theta_{\vec{\lambda}}$ on the calibration set of task $t$, and $M_t^{(l)}$ are the activations of the expert model $t$ on the same calibration data. Since searching a multi-dimensional grid on the full validation/test set is exponentially expensive, M-AOS offers a uniquely practical solution, performing the multi-dimensional search in seconds on a tiny calibration set of $N=128$ samples.

\section{Experimental Design}

In line with our empiricist philosophy, we designed a comprehensive and large-scale experimental suite to rigorously test AOS across a wide array of conditions. Our goal was to leave no stone unturned in validating our hypotheses and understanding the factors that influence the success of our method. The full study comprises over 2,160 unique experimental configurations.

\textbf{Tasks and Models.} We operate in a multi-task merging scenario. We begin with a ResNet-18 model \citep{he2016deep} pre-trained on CIFAR-100, which serves as our base model ($\theta_{base}$). We then create three specialist models by fine-tuning this base model for 2 epochs on three separate downstream image classification tasks: CIFAR-10 \citep{krizhevsky2009learning}, SVHN \citep{netzer2011reading}, and FashionMNIST \citep{xiao2017fashion}. Fine-tuning was performed using the AdamW optimizer \citep{loshchilov2017decoupled} with a learning rate of 2e-5. For each task, we merge the fine-tuned task model ($\theta_{task}$) with the original base model ($\theta_{base}$). The final reported accuracy for a given configuration is the average test accuracy achieved across these three independent merge experiments.

\textbf{Swept Parameters.} To ensure our conclusions are robust, we performed exhaustive sweeps over the key parameters of AOS:
\begin{itemize}
    \item \textbf{Number of Samples (N):} We tested the data efficiency of AOS by varying the number of samples used for activation calculation across the set $N \in \{8, 16, 32, 64, 128, 256\}$. These samples were drawn randomly from the training set of the respective downstream task.
    \item \textbf{Target Layer (L):} We investigated the effect of representation depth by extracting activations from three different blocks of the ResNet-18: the end of the third residual block (`layer3`), the end of the fourth residual block (`layer4`), and the output of the final global average pooling layer (`avgpool`).
    \item \textbf{Similarity Metric:} We compare the four metrics detailed in Section 3: CKA, MSE, Cosine, and MMD.
\end{itemize}

\textbf{Baselines and Oracle.} We compare AOS against two critical baselines:
\begin{itemize}
    \item \textbf{Weight Averaging (WA) Baseline:} This is a naive but common approach where models are simply averaged, corresponding to a fixed $\lambda=0.5$. It represents the performance one would achieve without any intelligent coefficient search.
    \item \textbf{Oracle Peak:} To establish the maximum possible performance for this merging strategy, we conducted a fine-grained grid search for $\lambda$ from 0.0 to 1.0 in steps of 0.01. For each of the 101 values, we constructed the merged model and evaluated it on the entire test set of the target task. The highest accuracy achieved in this search is reported as the "Oracle Peak." This is the performance ceiling that AOS aims to reach without the massive computational cost of the search.
\end{itemize}

\textbf{Reproducibility.} All experiments were repeated across 5 independent random seeds. The seeds control the initialization of the fine-tuning process, resulting in slightly different task models to be merged. We report the mean and standard deviation of the test accuracy across these 5 seeds for every configuration.

\begin{table*}[t]
\caption{Model merging classification accuracies (\%) across CIFAR-10, SVHN, and FashionMNIST under different AOS configurations. All results are averages and standard deviations across 5 independent seeds. WA refers to standard Weight Averaging backbone baseline. The Oracle Peak represents the best possible accuracy achievable by tuning $\lambda$ via a full validation set evaluation.}
\label{table:results}
\vskip 0.15in
\begin{center}
\begin{small}
\begin{sc}
\setlength{\tabcolsep}{2.5pt}
\begin{tabular}{lcccccc}
\toprule
\textbf{Config (N, Layer)} & \textbf{AOS-CKA} & \textbf{AOS-MSE} & \textbf{AOS-Cosine} & \textbf{AOS-MMD} & \textbf{Oracle Peak} & \textbf{WA Baseline} \\
\midrule
        N=256, L=avgpool & \textbf{58.43\% $\pm$ 0.00\%} & 56.83\% $\pm$ 0.49\% & 57.81\% $\pm$ 0.00\% & 54.63\% $\pm$ 0.00\% & 58.43\% $\pm$ 0.00\% & 34.58\% \\
        N=256, L=4 & \textbf{58.43\% $\pm$ 0.00\%} & 56.83\% $\pm$ 0.49\% & 57.81\% $\pm$ 0.00\% & 54.63\% $\pm$ 0.00\% & 58.43\% $\pm$ 0.00\% & 34.58\% \\
        N=256, L=3 & 58.35\% $\pm$ 0.07\% & 51.68\% $\pm$ 0.00\% & 54.63\% $\pm$ 0.00\% & 47.72\% $\pm$ 0.00\% & 58.43\% $\pm$ 0.00\% & 34.58\% \\
\midrule
        N=128, L=avgpool & 58.34\% $\pm$ 0.19\% & 56.83\% $\pm$ 0.49\% & 57.66\% $\pm$ 0.57\% & 54.63\% $\pm$ 0.00\% & 58.43\% $\pm$ 0.00\% & 34.58\% \\
        N=128, L=4 & 58.34\% $\pm$ 0.19\% & 56.83\% $\pm$ 0.49\% & 57.66\% $\pm$ 0.57\% & 54.63\% $\pm$ 0.00\% & 58.43\% $\pm$ 0.00\% & 34.58\% \\
        N=128, L=3 & 58.35\% $\pm$ 0.07\% & 51.68\% $\pm$ 0.00\% & 54.63\% $\pm$ 0.00\% & 48.51\% $\pm$ 1.58\% & 58.43\% $\pm$ 0.00\% & 34.58\% \\
\midrule
        N=64, L=avgpool & 58.31\% $\pm$ 0.18\% & 57.32\% $\pm$ 0.60\% & 57.66\% $\pm$ 0.57\% & 55.80\% $\pm$ 0.95\% & 58.43\% $\pm$ 0.00\% & 34.58\% \\
        N=64, L=4 & 58.31\% $\pm$ 0.18\% & 57.32\% $\pm$ 0.60\% & 57.66\% $\pm$ 0.57\% & 55.80\% $\pm$ 0.95\% & 58.43\% $\pm$ 0.00\% & 34.58\% \\
        N=64, L=3 & 58.35\% $\pm$ 0.07\% & 51.68\% $\pm$ 0.00\% & 54.63\% $\pm$ 0.00\% & 47.72\% $\pm$ 0.00\% & 58.43\% $\pm$ 0.00\% & 34.58\% \\
\midrule
        N=32, L=avgpool & 58.00\% $\pm$ 0.24\% & 57.32\% $\pm$ 0.60\% & 58.00\% $\pm$ 0.24\% & 55.27\% $\pm$ 1.27\% & 58.43\% $\pm$ 0.00\% & 34.58\% \\
        N=32, L=4 & 58.00\% $\pm$ 0.24\% & 57.32\% $\pm$ 0.60\% & 58.00\% $\pm$ 0.24\% & 55.27\% $\pm$ 1.27\% & 58.43\% $\pm$ 0.00\% & 34.58\% \\
        N=32, L=3 & 58.35\% $\pm$ 0.07\% & 52.27\% $\pm$ 1.18\% & 54.63\% $\pm$ 0.00\% & 49.30\% $\pm$ 1.94\% & 58.43\% $\pm$ 0.00\% & 34.58\% \\
\midrule
        N=16, L=avgpool & 57.68\% $\pm$ 0.73\% & 57.56\% $\pm$ 0.49\% & 58.00\% $\pm$ 0.24\% & 54.23\% $\pm$ 2.20\% & 58.43\% $\pm$ 0.00\% & 34.58\% \\
        N=16, L=4 & 57.68\% $\pm$ 0.73\% & 57.56\% $\pm$ 0.49\% & 58.00\% $\pm$ 0.24\% & 54.23\% $\pm$ 2.20\% & 58.43\% $\pm$ 0.00\% & 34.58\% \\
        N=16, L=3 & 58.22\% $\pm$ 0.21\% & 52.86\% $\pm$ 1.45\% & 55.41\% $\pm$ 0.95\% & 50.09\% $\pm$ 1.94\% & 58.43\% $\pm$ 0.00\% & 34.58\% \\
\midrule
        N=8, L=avgpool & 56.05\% $\pm$ 1.70\% & 57.07\% $\pm$ 0.60\% & 57.66\% $\pm$ 0.57\% & 55.07\% $\pm$ 2.08\% & 58.43\% $\pm$ 0.00\% & 34.58\% \\
        N=8, L=4 & 56.05\% $\pm$ 1.70\% & 57.07\% $\pm$ 0.60\% & 57.66\% $\pm$ 0.57\% & 55.07\% $\pm$ 2.08\% & 58.43\% $\pm$ 0.00\% & 34.58\% \\
        N=8, L=3 & 57.56\% $\pm$ 0.62\% & 52.27\% $\pm$ 1.18\% & 55.41\% $\pm$ 0.95\% & 50.69\% $\pm$ 2.65\% & 58.43\% $\pm$ 0.00\% & 34.58\% \\
\bottomrule
\end{tabular}
\end{sc}
\end{small}
\end{center}
\vskip -0.1in
\end{table*}

\section{Results}

The results of our exhaustive empirical study, summarized in Table \ref{table:results}, provide unequivocal support for Activation Overlap Search. The data demonstrates not only the effectiveness of our proposed method but also the stark failure of naive approaches and the clear superiority of certain configurations of AOS.

\textbf{AOS-CKA Perfectly Recovers Oracle Performance.} The most striking result is the performance of AOS with the CKA metric. At sufficient data sample sizes (N=256) and using activations from deeper layers (`avgpool` or `layer4`), AOS-CKA achieves an average accuracy of 58.43\%. This exactly matches the Oracle Peak performance, which was found by an expensive and computationally intensive grid search. This finding is the central validation of our core hypothesis: maximizing CKA similarity in activation space is a near-perfect proxy for maximizing final task accuracy.

\textbf{Baselines and the Need for a Principled Search.} The necessity of finding the correct $\lambda$ is made clear by the complete failure of the Weight Averaging (WA) baseline. With a fixed $\lambda=0.5$, the merged model's performance plummets to 34.58\%, a staggering 24-point drop from the oracle peak. This demonstrates that blindly averaging models that have diverged during fine-tuning can be catastrophic. AOS provides the principled search needed to avoid this failure mode.

\textbf{Remarkable Data Efficiency.} A key advantage of AOS is its practicality. Our results show that this high performance is not contingent on using a large dataset for activation calculation. With just N=16 samples, AOS-CKA (using `avgpool` or `layer4`) achieves 57.68\% accuracy, and AOS-Cosine reaches 58.00\%. These scores are within 1\% of the oracle peak and are over 23 points higher than the WA baseline. This data efficiency makes AOS a lightweight and practical tool for any model merging pipeline. Even at the extreme low-data regime of N=8, AOS-based methods significantly outperform the baseline, although CKA's performance begins to degrade due to the instability of estimating covariance matrices from few samples.

\textbf{The Superiority of CKA.} Across the full range of experiments, CKA consistently emerges as the most reliable and highest-performing similarity metric. While AOS-Cosine and AOS-MSE deliver respectable results, often scoring in the 56-58\% range, they are less consistent across layers and sample sizes and never fully close the gap to the oracle. AOS-MMD consistently performs the worst among the tested metrics, suggesting that comparing the full activation distributions with an RBF kernel is not as effective as comparing the geometric structure of the manifolds, which CKA excels at.

\textbf{The Importance of Batch Normalization Statistics.} It is critical to emphasize that none of these results, including the Oracle Peak, would be fully observable without our proposed fix for handling Batch Normalization running statistics. Initial experiments without this fix resulted in numerical instability and NaN values for any $\lambda$ above approximately 0.3. This completely prevented the oracle search from finding the true peak and would have unfairly handicapped all AOS methods. Our simple averaging of BN \texttt{running\_mean} and \texttt{running\_var} completely stabilized the process, enabling a fair and thorough empirical investigation. This is a crucial practical finding for the model merging community.

\subsection{Batch Normalization Merging Ablation Study}
To systematically understand the role of Batch Normalization (BN) running statistics in model merging, we conducted a rigorous ablation study comparing four distinct strategies:
\begin{enumerate}
    \item \textbf{Averaged BN Statistics (Ours):} The running mean and running variance buffers are computed as a simple average of the fine-tuned experts' buffers, which is a stable and mathematically coherent initialization.
    \item \textbf{Base Model BN Statistics:} The running statistics of the pre-trained base model are used directly without any merging or adaptation.
    \item \textbf{Naive BN Interpolation:} The running statistics are linearly interpolated using the merge coefficient $\lambda$ as: $\mu_{\lambda} = (1-\lambda)\mu_{\text{base}} + \lambda \mu_{\text{expert}}$ and $\sigma^2_{\lambda} = (1-\lambda)\sigma^2_{\text{base}} + \lambda \sigma^2_{\text{expert}}$.
    \item \textbf{BN Recalibration:} The running statistics of the merged model are fully re-estimated by performing a single forward pass of the $N=128$ calibration samples while setting the BN layers to training mode (with a momentum of 1.0 to capture exact batch statistics).
\end{enumerate}
All four strategies are evaluated using $N=128$, target layer \texttt{layer4}, and seed 42.

\begin{table}[h]
\caption{Ablation study on Batch Normalization merging strategies. AOS-CKA Lambda and Acc represent the selected coefficient and final test accuracy under our framework. Oracle Lambda and Acc show the absolute peak performance and its coefficient found via an exhaustive test-set grid search.}
\label{table:bn_ablation}
\vskip 0.15in
\begin{center}
\begin{small}
\begin{sc}
\setlength{\tabcolsep}{1.5pt}
\begin{tabular}{lcccc}
\toprule
\textbf{BN Strategy} & \textbf{AOS $\lambda$} & \textbf{AOS Acc.} & \textbf{Oracle $\lambda$} & \textbf{Oracle Acc.} \\
\midrule
Averaged & 0.75 & 58.43\% & 0.75 & 58.43\% \\
Base Model & 0.45 & 30.04\% & 0.60 & 31.82\% \\
Naive Interp. & 1.05 & 8.90\% & 0.90 & 55.71\% \\
Recalibrate & 0.80 & \textbf{69.78\%} & 1.05 & \textbf{71.14\%} \\
\bottomrule
\end{tabular}
\end{sc}
\end{small}
\end{center}
\vskip -0.1in
\end{table}

The empirical findings of this study, summarized in Table \ref{table:bn_ablation}, reveal several critical insights:
First, using the base model's unmerged BN statistics (\textbf{Base Model}) results in severe performance degradation, achieving a maximum Oracle accuracy of only 31.82\%. This is because the unadapted statistics fail to represent the massive shift in activation distributions introduced by task fine-tuning.
Second, \textbf{Naive Interpolation} of BN running statistics suffers from extreme numerical instability. As the merge coefficient $\lambda$ exceeds $1.00$, the interpolated variance becomes mathematically negative or poorly conditioned, causing PyTorch to output NaN values. This leads to immediate activation and accuracy collapse (8.90\%), empirically validating our initial bug diagnosis.
Third, \textbf{BN Recalibration} emerges as a spectacular enhancement. By re-estimating the mean and variance on a tiny calibration batch of 128 samples, the peak Oracle accuracy of the merged model leaps from 58.43\% to \textbf{71.14\%} (a massive \textbf{+12.71\%} absolute improvement). Crucially, AOS-CKA remains highly effective under this strategy, identifying $\lambda=0.80$ which achieves \textbf{69.78\%} test accuracy, recovering 98.1\% of the newly elevated Oracle Peak. This indicates that representation similarity in activation space remains a robust and high-fidelity proxy even when BN running statistics are adaptively recalibrated.

\subsection{Scaling to Multidimensional Merging}
To evaluate the capability of Multidimensional AOS (M-AOS), we perform sweeps over task-specific coefficients $(\lambda_{\text{cifar}}, \lambda_{\text{svhn}}, \lambda_{\text{fmnist}})$ in the range $[0.1, 1.1]$. We compare M-AOS against both standard Weight Averaging (WA) and the single-coefficient AOS-CKA framework. The results, evaluated across 5 independent seeds with $N=128$ and target layer \texttt{layer4}, are shown in Figure \ref{fig:maos_comparison}.

\begin{figure}[ht]
\vskip 0.2in
\begin{center}
\centerline{\includegraphics[width=\columnwidth]{maos_comparison.png}}
\caption{Average multi-task test accuracy (\%) across 5 independent seeds. M-AOS CKA (59.67\% $\pm$ 0.23\%) significantly outperforms single-coefficient AOS-CKA (58.34\% $\pm$ 0.19\%) and the WA baseline (34.58\% $\pm$ 0.00\%), closing the gap to the multi-dimensional Oracle Peak (60.08\% $\pm$ 0.06\%).}
\label{fig:maos_comparison}
\end{center}
\vskip -0.2in
\end{figure}

The empirical findings are highly encouraging:
\begin{enumerate}
    \item \textbf{Significant Performance Boost:} M-AOS CKA achieves an average test accuracy of \textbf{59.67\% $\pm$ 0.23\%}, outperforming the single-coefficient AOS-CKA ($58.34\% \pm 0.19\%$) by a substantial \textbf{+1.33\%} in absolute accuracy. This confirms that finding task-specific coefficients is highly beneficial for multi-task merging.
    \item \textbf{Near-Oracle Performance:} M-AOS CKA recovers the vast majority of the gains of the multidimensional search, coming within \textbf{0.41\%} of the absolute multidimensional Oracle Peak of \textbf{60.08\% $\pm$ 0.06\%}.
    \item \textbf{Alternative Metrics Performance:} While CKA is the top performer, M-AOS with Cosine similarity also performs strongly at 58.19\% $\pm$ 0.97\%, followed by MSE at 57.71\% $\pm$ 0.00\% and MMD at 56.10\% $\pm$ 1.21\%.
\end{enumerate}
This demonstrates that M-AOS is a highly robust, scalable, and practical solution for multi-task model merging, enabling efficient task-specific coefficient tuning without the prohibitive cost of multi-dimensional grid search on full validation sets.

\section{Discussion}

Our extensive results provide a clear picture of how to successfully apply AOS, but they also raise important questions about the underlying mechanics. We now discuss the implications of our findings.

\textbf{Ablation of Distance Metrics: Why CKA Excels.} The superior performance of CKA suggests that the key to a successful merge is not preserving the exact coordinate values of the activations (as MSE would measure) nor just their direction (as Cosine similarity measures). Instead, CKA's success points to the importance of preserving the relational structure of the activation manifold. CKA is invariant to rotation and uniform scaling, meaning it focuses on whether a set of inputs $\{x_i\}$ produce a similar geometry of activations $\{f(x_i)\}$. This geometric or relational similarity appears to be a much better proxy for preserved functional capabilities than metrics sensitive to the specific basis of the activation space. MSE is likely too strict, penalizing valid solutions that are simply rotated or shifted. MMD may be too general, as its focus on matching entire distributions might be less relevant than matching the instance-wise structure that CKA captures.

\textbf{Ablation of Layers: Deeper is Better.} The results consistently show that using activations from deeper layers (`layer4`, `avgpool`) yields better estimates of $\lambda$ than using a shallower layer (`layer3`), especially for metrics like MSE and Cosine. This aligns with our intuition that model merging should primarily aim to combine high-level semantic capabilities. Deeper layers in a CNN are known to capture more abstract, task-relevant features. Attempting to align representations at shallower layers, which may encode more generic, low-level features like edges and textures, might be an incorrect or overly restrictive objective. The final `avgpool` layer, which represents a summary of the image's features before classification, appears to be an excellent and reliable choice.

\textbf{The Centrality of BN Statistics in Merging.} Our discovery regarding the BN running statistics should be considered a key contribution. It highlights that a significant part of a model's learned function at inference time is encoded not just in its weights but also in these statistics. Any merging method that manipulates weights containing BN layers must also have a coherent strategy for combining their statistics. Our simple interpolation method works exceptionally well, but this opens up avenues for future work, such as exploring more sophisticated ways to merge these statistics, perhaps informed by the data distribution. We believe many reported failures or instabilities in model merging literature may be attributable to this often-overlooked component.

\textbf{Limitations and Future Work.} Our empirical study, while vast, is not exhaustive. We focused on a single architecture (ResNet-18) and vision tasks. A crucial next step is to validate the effectiveness of AOS on Transformer-based models in natural language processing, where model merging is also highly prevalent. While our initial formulation of AOS focused on finding a single, global scalar $\lambda$, we successfully extended this to Multidimensional AOS (M-AOS) to find task-specific coefficients. A natural next step is to explore layer-specific coefficients ($\lambda_l$) or even more fine-grained, block-specific ones. M-AOS could be adapted to find these coefficients in a hierarchical manner. Finally, investigating other, potentially more powerful, representation similarity metrics is a promising direction for future research.

\section{Conclusion}

In this paper, we introduced Activation Overlap Search (AOS), a data-driven, training-free, and computationally efficient method for estimating the interpolation coefficient in model merging. Grounded in a philosophy of rigorous empirical validation, we conducted a large-scale study that spanned thousands of unique configurations. Our results are definitive: AOS, when paired with the Centered Kernel Alignment (CKA) similarity metric, serves as a near-perfect proxy for the expensive process of evaluating a merged model's performance on a validation set. It successfully identified the oracle peak accuracy of 58.43\% in our testbed, a feat that naive Weight Averaging (34.58\%) failed at catastrophically. We demonstrated the method's extreme data efficiency, achieving near-optimal results with as few as 16 data samples. We also identified and provided a crucial fix for handling Batch Normalization statistics, a key practical challenge in model merging. Furthermore, we extended our framework to Multidimensional AOS (M-AOS), which automates the search for task-specific merging coefficients. Under a multidimensional grid, M-AOS with CKA achieved an average accuracy of 59.67\% across 5 independent seeds, significantly outperforming single-coefficient methods and closing the gap to the multidimensional Oracle Peak (60.08\%) to within 0.41\%.

Our work champions a functional-space perspective on model merging, shifting the focus from the parameter space to the more direct and informative activation space. By replacing computationally prohibitive grid searches with a principled, data-driven proxy, AOS offers a practical and powerful tool that enhances the viability and effectiveness of model merging techniques.

% In the unusual situation where you want a paper to appear in the
% references without citing it in the main text, use \nocite
\nocite{langley00}

\bibliography{paper}
\bibliographystyle{icml2026}

\end{document}